\documentclass[book.tex]{subfiles}

\begin{document}

\chapter{Spin and Fermionic Path Integrals}

\label{chap:spin_fermionic}
\noindent\rule{11cm}{0.4pt} \\
\textbf{Prerequisites:}  \autoref{chap:path_integrals} \\
\textbf{Difficulty Level:} ** \\
\noindent\rule{11cm}{0.4pt}
\vspace{5mm}

\textcolor{red}{TODO: Cite Shankar}

Suppose we have a spin system with $S$ degrees of freedom. We can model this spin with a $2S+1$ dimensional Hilbert space $\mathcal{H}$. We can then resolve the identity as
\begin{align}
    I = \sum_{-S}^{S}|S_Z\rangle \langle S_Z |
\end{align}

\textcolor{red}{TODO: I'm having trouble following the discussion in Shankar. I think the idea is to use a non-standard resolution of identity. We eventually reach the propagator definition here}

\begin{align}
\langle \Omega_f | U(t) | \Omega_i \rangle &= \int [\mathcal{D}\Omega] \exp \Bigg[i\int_{t_1}^{t_2} [S\cos \theta \dot{\phi} - \mathcal{H}(\Omega)] dt\Bigg]
\end{align}

(TODO: Shankar notes there are apparently major convergence issues here)

\section{Fermionic Oscillator}



\section{Fermionic Path Integrals}



\end{document}
